\documentclass[11pt]{article}
\usepackage[margin=1in]{geometry}
\usepackage[T1]{fontenc}
\usepackage{lmodern,amsmath,amssymb,amsthm,mathtools}
\usepackage[colorlinks=true,allcolors=blue]{hyperref}
\newtheorem{theorem}{Theorem}
\newtheorem{lemma}[theorem]{Lemma}
\newtheorem{corollary}[theorem]{Corollary}
\newtheorem{proposition}[theorem]{Proposition}
\theoremstyle{remark}
\newtheorem{remark}[theorem]{Remark}
\newcommand{\R}{\mathbb R}
\newcommand{\PP}{\mathcal P}
\newcommand{\dd}{\,\mathrm d}
\newcommand{\norm}[1]{\left\lVert#1\right\rVert}
\newcommand{\curl}{\operatorname{curl}}
\newcommand{\diver}{\operatorname{div}}
\title{A history criterion for the non-axisymmetric part of a viscous flow}
\author{}
\date{September 27, 2026}

\begin{document}
\maketitle
\begin{abstract}
We give a sufficient regularity criterion when the angular mean of a
three-dimensional Navier--Stokes velocity satisfies prescribed anisotropic
bounds. The remaining assumptions are two time integrals of the angular
fluctuation. A radial circulation barrier preserves a spatial factor in the
Reynolds torque, while a divergence representation of the Reynolds stress
avoids differentiating the fluctuation in the ancient-limit argument.
For finite-energy whole-space flows, one weighted amplitude integral
suffices: the exact defect equation supplies the missing gradient history
on short intervals following the mean flow, with the pressure tail retained.
The proof adapts the axisymmetric argument of Constantin, Ignatova, and
Vicol. This is an analytic conditional criterion, not a Lean theorem or
an arbitrary-data regularity proof.
\end{abstract}

\section{Statement}

Let $Q=B(1)\times(-1,0)$ and let $(u,p)$ be a suitable weak solution of
\[
 u_t+(u\cdot\nabla)u-\Delta u+\nabla p=0,\qquad \diver u=0,
\]
smooth on compact subsets of $Q$. Assume the usual local energy class
$u\in L^\infty_tL^2_x\cap L^2_tH^1_x$ and $p\in L^{3/2}(Q)$.
For rotations $Q_\phi$ about the fixed third axis put
\[
 \PP u(x)=\frac1{2\pi}\int_0^{2\pi}Q_\phi u(Q_\phi^{-1}x)\dd\phi,
 \qquad v=\PP u,\qquad w=u-v.
\]
The same notation denotes the induced action on tensors. Write $r=|x'|$,
$z=x_3$, and $\tau=-t$. Unless a smaller ball is specified, suprema are on $B(1)$.

\begin{theorem}\label{thm:main}
Suppose $0<h<1/2$ and, for every $a+b\leq2$,
\begin{align}
 |\partial_r^a\partial_z^b v_r|
 &\leq C\tau^{-1/2-a/2-(1/2-h)b},\label{eq:meanr}\\
 |\partial_r^a\partial_z^b v_z|+|\partial_r^a\partial_z^b v_\theta|
 &\leq C\tau^{-1/2-h-a/2-(1/2-h)b}.\label{eq:meanz}
\end{align}
Set $K(t)=\norm{w(t)}_\infty$ and
$\mathcal T(t)=\norm{(\PP[(w\cdot\nabla)w])_\theta}_\infty$.
If
\begin{equation}\label{eq:history}
 \int_{-1}^0 K(t)^2\dd t<\infty,
 \qquad
 \int_{-1}^0\sqrt{-t}\,\mathcal T(t)\dd t<\infty,
\end{equation}
then $(0,0)$ is a regular point of $u$.
\end{theorem}

Angular averaging in $\mathcal T$ precedes the absolute value.
The mean bounds are those used in \cite{CIV}. Their validity and
\eqref{eq:history} are hypotheses, not consequences of the energy inequality.
The axis and the ball are fixed; the theorem makes no assertion about
propagating approximate symmetry from a shrinking core.

\begin{corollary}\label{cor:powers}
In Theorem~\ref{thm:main}, a sufficient replacement for
\eqref{eq:history}, with $L(t)=\norm{\nabla w(t)}_\infty$ using
Cartesian derivatives, is
\[
 K(t)\leq C\tau^{-\beta},\qquad
 L(t)\leq C\tau^{-\beta-\sigma},\qquad
 \beta<\tfrac12,\quad 2\beta+\sigma<\tfrac32,
\]
with $\beta,\sigma\geq0$. In particular, the derivative cost
$\sigma=1/2$ permits every $\beta<1/2$.
\end{corollary}

Indeed $\mathcal T\leq KL$, so a sufficient history condition is
$\int K^2<\infty$ and $\int\sqrt{-t}\,KL<\infty$.
These are two separate bounds. No derivative estimate is inferred from
the amplitude estimate in this corollary. Section~\ref{sec:smoothing}
gives such an implication under additional whole-space hypotheses.
Another sufficient choice near zero is
$K\leq C\tau^{-1/2}\log(e/\tau)^{-p}$ and
$L\leq C\tau^{-1}\log(e/\tau)^{-q}$, where $p>1/2$ and $p+q>1$.

\section{Two scalar history estimates}

\begin{lemma}\label{lem:hardy}
Let $f\geq0$ be continuous on compact subsets of $(0,a]$, with
$J_f=\int_0^a\sqrt{s}f(s)\dd s<\infty$, and define
$B(\tau)=\int_\tau^a f(s)\dd s$. Then
\begin{equation}\label{eq:scalar}
 \sqrt\tau B(\tau)\longrightarrow0,
 \qquad \int_0^a s^{-1/2}B(s)\dd s=2J_f.
\end{equation}
For $0\leq\eta<1/2$,
\begin{equation}\label{eq:hardy}
 \norm{\tau^{-\eta}\int_\tau^a s^\eta f(s)\dd s}_{L^2(0,a)}
 \leq\frac{J_f}{\sqrt{1-2\eta}}.
\end{equation}
\end{lemma}
\begin{proof}
For the limit, apply dominated convergence to
$\mathbf1_{s\geq\tau}\sqrt{\tau/s}\,\sqrt{s}f(s)$.
The equality is Tonelli's theorem. Minkowski's integral inequality gives
\[
 \norm{\int_0^a\mathbf1_{\tau<s}\tau^{-\eta}s^\eta f(s)\dd s}_2
 \leq\int_0^a s^\eta f(s)
       \left(\int_0^s\tau^{-2\eta}\dd\tau\right)^{1/2}\dd s,
\]
which is \eqref{eq:hardy}. Truncation justifies the calculation before
passing to the possibly singular endpoint.
\end{proof}

\begin{lemma}\label{lem:barrier}
Suppose an axisymmetric circulation $\Gamma=r v_\theta$ satisfies
\[
 \Gamma_t+v_r\Gamma_r+v_z\Gamma_z
   =(\partial_{rr}-r^{-1}\partial_r+\partial_{zz})\Gamma+rF_\theta,
\]
on a ball for $-a<t<0$, with bounded initial and lateral values,
$\Gamma=0$ on the axis, $|v_r|\leq C_r(-t)^{-1/2}$, and
$|F_\theta|\leq f(-t)$ as in Lemma~\ref{lem:hardy}. Then
\begin{equation}\label{eq:barrier}
 |\Gamma(r,z,t)|\leq A(-t)+rB(-t),\quad
 A(\tau)=C_0+C_r\int_\tau^a s^{-1/2}B(s)\dd s,
\end{equation}
where $C_0$ bounds the initial and lateral values. In particular $A$
is bounded, and $\sqrt\tau B(\tau)\to0$.
\end{lemma}
\begin{proof}
For $r>0$, applying the circulation operator to
$\Phi(r,t)=A(-t)+rB(-t)$ gives
\[
 C_r(-t)^{-1/2}B+r f+v_rB+B/r\geq r f.
\]
Thus $\Phi$ dominates both $\Gamma$ and $-\Gamma$ by the maximum
principle on each compact time prefix. A positive maximum of their
difference cannot occur on the axis, where $\Gamma=0$ and $\Phi\geq0$;
no differentiability of $r$ at the axis is required. The bounds follow
from Lemma~\ref{lem:hardy}.
\end{proof}

\section{The mean equation and its rescaled stress}

Rotation covariance and $\PP w=0$ give the exact equation
\begin{equation}\label{eq:mean}
 v_t+(v\cdot\nabla)v-\Delta v+\nabla\PP p=F,
 \qquad F=-\PP[(w\cdot\nabla)w]=-\diver\PP(w\otimes w).
\end{equation}
In particular $\norm{F(t)}_\infty\leq K(t)L(t)$, using the Euclidean
operator norm for $\nabla w$. Other tensor norm conventions change only
a fixed constant.

The original solution has a regular spherical annulus inside any
prescribed interval of radii. Indeed, partial regularity makes the
top-slice singular set have zero $\mathcal H^1$ measure. Its image under
$x\mapsto|x|$ has measure zero, so some intervening sphere is regular.
Compactness and interior estimates give an annulus on which $u$,
$\nabla u$, and $\nabla^2u$ are bounded through the endpoint, as in the
first part of \cite[Lemma~3.2]{CIV}. Its rotational invariance gives the
same bounds for $v$ and $w$. This argument uses suitability of $u$;
it does not require a partial-regularity theorem for \eqref{eq:mean}.

Choose a sphere in this annulus, and take $f(\tau)$ to be the supremum
of $|F_\theta(-\tau)|$ on its enclosed closed ball. This is continuous on
compact time prefixes and satisfies $J_f<\infty$ by
\eqref{eq:history}. Lemma~\ref{lem:barrier} applies to the mean.

Write the meridional entries of $\PP(w\otimes w)$ in the cylindrical
frame as $a=R_{rr}$, $b=R_{\theta\theta}$, $c=R_{zz}$, and $d=R_{rz}$.
Their absolute values are at most $K^2$. Direct differentiation yields
\begin{align}
 F_r&=-a_r-(a-b)/r-d_z,& F_z&=-d_r-d/r-c_z,\label{eq:force}\\
 \frac{(\curl F)_\theta}{r}
 &=\left(\partial_{rr}+\frac3r\partial_r-\partial_{zz}\right)\frac d r
   +\frac1r\partial_{rz}(c-a)-\frac1{r^2}\partial_z(a-b).
 \label{eq:lift}
\end{align}
The quantities here are smooth at the axis in the appropriate lifted
variables. For a distributional estimate it suffices that $d=O(r)$.

Consider the rescaling used in \cite[Section~4]{CIV}: for $t_n\uparrow0$,
\[
 \tau_n=-t_n,\quad \lambda_n=\tau_n^{1/2},\quad
 \delta_n=\tau_n^h,\quad \mu_n=\lambda_n/\delta_n,\quad
 r=\lambda_n R,\quad z=z_n+\mu_n Z,\quad t=\tau_n s.
\]
The mean components are normalized as
$V_n=\lambda_n v_r$, $W_n=\lambda_n\delta_n v_z$,
$S_n=\lambda_n\delta_n v_\theta$, and the potential vorticity as
$\Omega_n=\lambda_n^3\delta_n(\curl v)_\theta/r$.
For a stress component $a$, let $a_n$ denote its composition with this
change of variables, without an amplitude normalization.

Lift $R=|X|$ to four radial dimensions. Since
\[
 \frac1R\partial_R H
   =\diver_X\left(\frac{XH}{R^2}\right)-\frac{2H}{R^2},
\]
the rescaled forcing in the potential-vorticity equation is
\begin{equation}\label{eq:scaled}
\begin{split}
 G_n=\lambda_n^2\delta_n\bigg[&
 (\Delta_X-\delta_n^2\partial_{ZZ})(d_n/R)\\
 &+\delta_n\partial_Z\left\{
  \diver_X\frac{X(c_n-a_n)}{R^2}
  -\frac{2(c_n-a_n)+(a_n-b_n)}{R^2}\right\}\bigg].
\end{split}
\end{equation}
Both $R^{-1}$ and $R^{-2}$ are locally integrable in $\R^4$.
The divergence identity has no axis-supported mass: the flux through
$R=\epsilon$ is $O(\epsilon^2)$ for bounded $H$.
Move all derivatives in \eqref{eq:scaled} onto a compactly supported
smooth test function with time support $J\Subset(-\infty,0)$.
The resulting absolute value is bounded by
\begin{equation}\label{eq:weaksmall}
 C_J\lambda_n^2\delta_n\int_J K(\tau_n s)^2\dd s
 =C_J\delta_n\int_{\tau_nJ}K(t)^2\dd t\longrightarrow0.
\end{equation}
Thus no derivative of the fluctuation is needed in this part of the
limiting argument. In a recentered zoom with $r=\lambda_n(A_n+R)$ and
$A_n\to\infty$, the azimuthal-vorticity force has the same prefactor
$\lambda_n^2\delta_n$ multiplying at most two derivatives of the stress;
its additional coefficients $(A_n+R)^{-1}$ and $(A_n+R)^{-2}$ are
uniformly bounded on compact sets. Estimate~\eqref{eq:weaksmall}
therefore applies there too.

\section{Meridional smallness}

Define, on a smaller ball,
\[
 G_v(t)=\norm{|\partial_r v_r|+|v_r/r|
                    +|\partial_z v_z|+|\partial_z v_r|}_\infty.
\]
We verify the two changes needed in the proof of
\cite[Proposition~1.5]{CIV} to conclude
\begin{equation}\label{eq:Gsmall}
 (-t)G_v(t)\longrightarrow0.
\end{equation}
First, the circulation barrier gives, in the finite-axis case,
\[
 |R S_n|\leq\delta_n A(\tau_n|s|)
             +\delta_n\lambda_n R B(\tau_n|s|)\longrightarrow0
\]
uniformly on compact rescaled time slabs and bounded $R$. The mean
bounds also give $|S_n/R|\leq C$ on those slabs, including the axis.
Splitting at $R=\epsilon$ shows that $S_n^2/R^2$ tends to zero locally
in $L^1$ in the lifted space. In the receding-axis case use instead
\[
 |S_n|\leq\frac{\delta_n A(\tau_n|s|)}{A_n+R}
                +\delta_n\lambda_n B(\tau_n|s|).
\]

Second, the force is small in a negative Sobolev norm integrated in
time, rather than uniformly small as a function. Away from the axis,
\eqref{eq:scaled} and \eqref{eq:weaksmall} give
\[
 \norm{G_n}_{L^1(J;W^{-2,p}(D))}\longrightarrow0
\]
for a fixed rectangle $D$ and any finite $p$. The analogous assertion
holds in the receding-axis case. The other terms in the rescaled
vorticity equation are uniformly bounded in this negative norm.
Together with the uniform spatial Lipschitz bound supplied by
\eqref{eq:meanr}--\eqref{eq:meanz}, this gives an increment bound
$C|s-s'|+\epsilon_n$ with $\epsilon_n\to0$. Interpolation gives
asymptotic uniform time equicontinuity, also at the selected upper
time. Finitely many initial members are smooth, so Arzel\`a--Ascoli
still applies. Boundedness and the small volume of $R<\epsilon$ give
convergence across the axis.

The remaining ancient-limit argument is unchanged: the force and swirl
source vanish, and the mean bounds supply the backward decay and
bounded Lipschitz drift required by \cite[Lemma~3.3]{CIV}. The limiting
vorticity vanishes. At the selected time, incompressibility and
boundedness of the limiting meridional velocity then contradict a
nonzero selected derivative, in either axis case. This proves
\eqref{eq:Gsmall}. No bound on $\partial_r v_z$ is asserted.

\section{Full-velocity enstrophy}

Choose a radial cutoff $\chi$ whose derivative is supported in a regular
annulus, and write
\[
 Y=\norm{\chi\curl u}_2^2,\quad
 M=\norm{\chi\nabla\curl u}_2^2,\quad
 V_\theta=\norm{v_\theta}_\infty.
\]
Here $V_\theta$ is taken on a ball containing the support of $\chi$,
with its boundary in the regular annulus; $K$ is still taken on $B(1)$.
All cutoff errors are uniformly bounded. Weighted $L^2$ contraction
of angular averaging, including its tensor action on gradients, gives
$Y_v\leq Y$ and $M_v\leq M$. Local div--curl estimates consequently
bound the first derivatives of $u,v$ by $C(Y+1)^{1/2}$ and their second
derivatives by $C(M+1)^{1/2}$ in these weighted norms.

For $\Pi_\chi(q)=\int\chi^2(\curl q\cdot\nabla)q\cdot\curl q$, put
$\omega_q=\curl q$. There is an exact telescoping identity
\begin{equation}\label{eq:telescope}
\begin{split}
 \Pi_\chi(u)-\Pi_\chi(v)=\int\chi^2\big[&
 (\omega_w\cdot\nabla)u\cdot\omega_u
 +(\omega_v\cdot\nabla)w\cdot\omega_u\\
 &+(\omega_v\cdot\nabla)v\cdot\omega_w\big].
\end{split}
\end{equation}
Integrate the derivative on $w$ off that factor in each term. The
remaining factors have derivative orders one and two, apart from
bounded cutoff errors. Hence, for every $\epsilon>0$,
\begin{align*}
 |\Pi_\chi(u)-\Pi_\chi(v)|
 &\leq CK(Y+1)^{1/2}(M+1)^{1/2}+C\\
 &\leq\epsilon M+C_\epsilon(K^2+1)(Y+1).
\end{align*}
For the mean's stretching, the cylindrical calculation in
\cite[Section~5]{CIV} gives
\[
 \Pi_\chi(v)\leq\epsilon M+C_\epsilon(G_v+V_\theta^2+1)(Y+1).
\]
In particular its large $\partial_r v_z$ term is integrated by parts
using $(\omega_v)_r=-\partial_z v_\theta$, with $v_\theta$ left
undifferentiated. Testing the original unforced vorticity equation
and absorbing both small dissipation terms gives
\begin{equation}\label{eq:energy}
 Y'+cM\leq C_*[G_v+V_\theta^2+K^2+1](Y+1).
\end{equation}

Fix $0<\eta<1/2$ and shorten the final interval so that
$|v_r/r|\leq\eta/\tau$, using \eqref{eq:Gsmall}. The equation for
$\tau^\eta v_\theta$ has nonnegative zeroth-order coefficient.
Its maximum principle therefore yields
\[
 V_\theta(-\tau)\leq C_\eta\tau^{-\eta}
                  +\tau^{-\eta}\int_\tau^a s^\eta f(s)\dd s.
\]
Lemma~\ref{lem:hardy} proves $V_\theta\in L^2_t$. Since $K\in L^2_t$
as well, Gronwall in \eqref{eq:energy} gives
\[
 Y(t)+1\leq C_\eta(-t)^{-C_*\eta}.
\]
Here $C_*$ was fixed before $\eta$. Choose $\eta$ with $2C_*\eta<1$.
Local Sobolev now places $u$ in $L^4_tL^6_x$ on a smaller cylinder.
The local criterion of \cite[Theorem~1.1(i)]{GKT}, applied on backward
cylinders ending at zero, makes $(0,0)$ regular. This finishes the
proof of Theorem~\ref{thm:main}.

\section{Paying the gradient history from the defect equation}
\label{sec:smoothing}

We now use the actual evolution of the defect to replace both histories
by one amplitude integral. The following whole-space assumptions let us
control the pressure produced by localization.

\begin{theorem}\label{thm:amplitude}
Suppose the solution in Theorem~\ref{thm:main} is the restriction of a
smooth solution on $\R^3\times(-1,0)$ with
$u\in L^\infty((-1,0);L^2(\R^3))$ and canonical pressure.
Retain the local suitability assumptions and the mean bounds
\eqref{eq:meanr}--\eqref{eq:meanz}. It suffices to assume, for some
$0<a<1$,
\begin{equation}\label{eq:amplitude}
 \int_0^a \tau^{-h/2}\sqrt{\log(e/\tau)}\,
               \norm{w(-\tau)}_{L^\infty(B(1))}^2\dd\tau<\infty.
\end{equation}
In particular, the amplitude estimate $\norm{w(-\tau)}_\infty
\leq C\tau^{-\beta}$ suffices if $\beta<1/2-h/4$, without a
separate derivative estimate. At equality, the additional factor
$\log(e/\tau)^{-p}$ suffices for $p>3/4$.
\end{theorem}

The theorem does not assert \eqref{eq:amplitude} from initial data.
Its proof has a whole-space smoothing estimate and a localization step.
Write $\mathbb L$ for the Leray projection, to distinguish it from
angular averaging $\PP$.

\begin{lemma}\label{lem:window}
Let a smooth divergence-free $U$ on $\R^3\times[0,d]$ satisfy
\[
 U_t-\Delta U=-\mathbb L\diver(U\otimes U)+g,
 \qquad V=\PP U,\quad W=U-V,\quad G=(1-\PP)g.
\]
Assume the fields and the spatial derivatives used below are bounded on
compact time prefixes, and that
$\norm V_\infty\leq M$, $\norm{\nabla V}_\infty\leq N$,
$\norm G_\infty\leq G_0$ there, with $M,N>0$.
Put $K(q)=\norm{W(q)}_\infty$, $L(q)=\norm{\nabla W(q)}_\infty$, and
\[
 Z_d=\sup_{0<q\leq d}\{K(q)+\sqrt q L(q)\},
 \qquad \log_+x=\max(\log x,0).
\]
There are universal constants such that
\begin{equation}\label{eq:window}
 Z_d\leq C_0(K(0)+dG_0)
 +C_1Nd\left[1+\log_+\frac{M}{N\sqrt d}\right]Z_d
 +C_2\sqrt d Z_d^2.
\end{equation}
Consequently $Z_d\leq2C_0(K(0)+dG_0)$ if the linear coefficient
is at most $1/4$ and $\sqrt d(K(0)+dG_0)$ is sufficiently small.
\end{lemma}

\begin{proof}
Rotation covariance gives the exact homogeneous-in-$W$ equation
\begin{equation}\label{eq:defect}
 W_t-\Delta W=-\mathbb L\diver\big[V\otimes W+W\otimes V
                          +(1-\PP)(W\otimes W)\big]+G.
\end{equation}
In particular the complete defect self-interaction is retained.
For the heat kernel $H_s$, the kernel of $e^{s\Delta}\mathbb L\diver$ is
\[
 (B_s)_{ijk}=\delta_{ij}\partial_k H_s
                         +\int_s^\infty\partial_{ijk}H_b\dd b.
\]
Differentiating the Gaussian and integrating in $b$ gives
\begin{align*}
 |B_s(y)|&\leq C(\sqrt s+|y|)^{-4},&
 |\nabla B_s(y)|&\leq C(\sqrt s+|y|)^{-5},\\
 \norm{B_s}_1&\leq C s^{-1/2},&
 \int |y|\,|\nabla B_s(y)|\dd y&\leq C s^{-1/2}.
\end{align*}
The first absolute moment of $B_s$ is not integrable at infinity.
Splitting the radial integral at $\sqrt s$ and $M/N$ instead gives
\begin{align}
 \int |B_s(y)|\min(N|y|,2M)\dd y
 &\leq CN\left[1+\log_+\frac{M}{N\sqrt s}\right],\label{eq:oseenlog}\\
 \int |\nabla B_s(y)|\min(N|y|,2M)\dd y
 &\leq CNs^{-1/2}.\label{eq:oseengrad}
\end{align}

For each endpoint $(x,q)$ follow the backward characteristic
$z'(s)=V(z(s),s)$, $z(q)=x$, and set $a(s)=V(z(s),s)$.
Move $a(s)\cdot\nabla W$ to the left side of \eqref{eq:defect};
its two linear tensors now contain $V-a$ instead of $V$, since
$\diver W=0$. This identity holds everywhere with this one
spatially constant $a(s)$. The translated Duhamel kernels at the
chosen endpoint are centered at $z(s)$, where
\[
 |V(y,s)-a(s)|\leq\min(N|y-z(s)|,2M).
\]
For the gradient estimate, differentiate that representation with
$a(s)$ held fixed. No derivative of the selected characteristic is
introduced: the same representation is valid at every nearby endpoint
with this fixed $a$. Estimates \eqref{eq:oseenlog}--\eqref{eq:oseengrad}
bound the linear amplitude and gradient terms.

For the quadratic gradient term move the derivative onto the tensor:
$\norm{\nabla(1-\PP)(W\otimes W)}_\infty\leq CKL$.
Angular averaging and its tensor action are contractions in the
Euclidean sup norms. The time integral
$\int_0^q(q-s)^{-1/2}s^{-1/2}\dd s=\pi$ then bounds this term
by $CZ_d^2$, before multiplication by $\sqrt q$.
The quadratic amplitude term is bounded by $C\sqrt q Z_d^2$.
The initial term and the divergence-free source use the ordinary
translated heat kernel. Finally
\[
 \int_0^q\left[1+\log_+\frac{M}{N\sqrt s}\right]\dd s
 \leq Cq\left[1+\log_+\frac{M}{N\sqrt q}\right].
\]
This proves \eqref{eq:window}. The estimate holds on every shorter
prefix, whose linear coefficient is no larger. A continuity bootstrap
gives the final claim; smoothness supplies the initial continuity.
If $K(0)+dG_0=0$, use uniqueness or a positive bootstrap bound tending
to zero.
\end{proof}

\begin{proof}[Proof of Theorem~\ref{thm:amplitude}]
First suppose a globally bounded defect equation as in
Lemma~\ref{lem:window} has been obtained, with $G$ uniformly bounded and
\begin{equation}\label{eq:meanlip}
 \norm{V(-\tau)}_\infty\leq C_M\tau^{-1/2-h},\qquad
 \norm{\nabla V(-\tau)}_\infty\leq C_N\tau^{-1-h}.
\end{equation}
Let $\ell(\tau)=\log(e/\tau)$ and choose
\begin{equation}\label{eq:dwindow}
 d(\tau)=\frac{c\tau^{1+h}}{\ell(\tau)},\qquad 0<c<1.
\end{equation}
On $[-\tau-d,-\tau]$, the linear coefficient in
\eqref{eq:window} is at most $C c(1+|\log c|+h)$: here
$Nd\leq C_Nc/\ell$ and $M/(N\sqrt d)$ is bounded by a
constant times $c^{-1/2}\tau^{-h/2}\sqrt\ell$.
Take $c$ small enough in terms of the mean bounds.

Write $K=\norm W_\infty$, $L=\norm{\nabla W}_\infty$, and
\[
 E(\tau)=\int_\tau^{\tau+d(\tau)} K(-s)^2\dd s.
\]
Condition \eqref{eq:amplitude}, with the localized global norm,
implies $K\in L^2_t$ and $E(\tau)\to0$.
There is a time $s_0\in[-\tau-d,-\tau-d/2]$ at which
$K(s_0)^2\leq2E(\tau)/d$.
Applying Lemma~\ref{lem:window} on $[s_0,-\tau]$ is therefore
legitimate for all sufficiently small $\tau$; its small parameter is
at most $C\sqrt{E(\tau)}+d^{3/2}G_0\to0$.
We obtain
\begin{equation}\label{eq:smoothing}
 K(-\tau)+\sqrt d L(-\tau)
 \leq C\big[d^{-1/2}\sqrt{E(\tau)}+dG_0\big].
\end{equation}
Consequently
\begin{equation}\label{eq:product}
 \sqrt\tau K(-\tau)L(-\tau)
 \leq C\sqrt\tau d^{-3/2}E(\tau)
                         +C\sqrt\tau d^{3/2}G_0^2.
\end{equation}
The second term is integrable. For the first, Tonelli gives
\begin{equation}\label{eq:tonelliwindow}
 \int_0^{a/2}\sqrt\tau d(\tau)^{-3/2}E(\tau)\dd\tau
 \leq C\int_0^a s^{-h/2}\sqrt{\ell(s)}K(-s)^2\dd s.
\end{equation}
Indeed, $\tau\leq s\leq\tau+d(\tau)$ implies
$s/(1+c)\leq\tau\leq s$. Since $d$ increases, these $\tau$ lie
in an interval of length at most $d(s)$, while the integrand's
coefficient is at most $Cs^{-1-3h/2}\ell(s)^{3/2}$.
This proves \eqref{eq:tonelliwindow}, hence the two histories on a
final interval, using $\mathcal T\leq KL$.

It remains to construct the localization used above. Choose a radial
cutoff $\chi$, equal to one on an interior ball and with its transition
inside a regular annulus of the original solution. Put
\[
 \Psi=(-\Delta)^{-1}\curl u,\qquad
 U=\curl(\chi\Psi)=\chi u+b,\qquad b=\nabla\chi\times\Psi.
\]
The $L^2$ hypothesis and incompressibility give $\curl\Psi=u$;
the potential is well defined by its order-minus-one kernel.
On the annulus, its separated kernel and every required derivative
are bounded against $\norm u_2$, while its local part is controlled
by interior spatial regularity. For canonical pressure, the same
decomposition for its order-zero transform of $u\otimes u$ bounds
$p$ and its spatial derivatives on the annulus. Likewise
\[
 \Psi_t=-\curl u-(-\Delta)^{-1}\curl\diver(u\otimes u)
\]
has bounded spatial derivatives there: its bilinear operator is of
order zero, with separated kernels controlled by $\norm u_2^2$.
These statements use only annular regularity and the uniform energy
bound, not regularity of the interior through time zero.

With pressure $\chi p$, the exact force of $U$ is
\begin{equation}\label{eq:cutoffforce}
\begin{split}
 f_U={}&b_t-2\nabla\chi\cdot\nabla u-(\Delta\chi)u-\Delta b
                       +p\nabla\chi\\
      &+(U\cdot\nabla)U-\chi(u\cdot\nabla)u.
\end{split}
\end{equation}
It is supported in the regular annulus and uniformly bounded in
spatial $C^1$. Its Leray projection is uniformly bounded, as follows
by subtracting $f_U(x)$ in the local principal value and using its
$L^1$ norm for the far part. Thus the equation for $U$ has
$g=\mathbb L f_U$ and bounded $G=(1-\PP)g$, including pressure.

Rotations preserve the interior ball, so $U=u$, $\PP U=v$, and
$U-\PP U=w$ there. The transition region contributes bounded fields
and derivatives; outside the cutoff $U=0$. The mean assumptions
therefore imply \eqref{eq:meanlip} globally. The Cartesian gradient
terms $v_r/r$ and $v_\theta/r$ are controlled by radial derivatives
along the segment to the axis. The global weighted amplitude integral
is bounded by \eqref{eq:amplitude} plus the finite contribution of a
bounded function. Equations \eqref{eq:smoothing}--\eqref{eq:tonelliwindow}
now supply the required histories on a smaller cylinder of the
original flow. Theorem~\ref{thm:main} applies there. The power and
logarithmic consequences follow by integrating
$\tau^{-h/2-2\beta}\ell(\tau)^{1/2-2p}$.
\end{proof}

\section{The sign of kinetic mean--defect transfer}

For a smooth finite-energy whole-space unforced flow, orthogonality of
angular averaging and integration by parts give
\begin{equation}\label{eq:kinetictransfer}
 \frac12\frac{\mathrm d}{\mathrm dt}\norm w_2^2
       +\nu\norm{\nabla w}_2^2=-T,
 \qquad
 \frac12\frac{\mathrm d}{\mathrm dt}\norm v_2^2
       +\nu\norm{\nabla v}_2^2=T,
 \quad T=\int w\cdot(w\cdot\nabla)v.
\end{equation}
We display viscosity explicitly in this section. Defect self-advection
and pressure do no kinetic work on $w$; the averaging correction is
orthogonal to $w$. In particular the signed net kinetic transfer obeys
\[
 -\int_0^t T(s)\dd s
 =\tfrac12\norm{v(0)}_2^2-\tfrac12\norm{v(t)}_2^2
                         -\nu\int_0^t\norm{\nabla v}_2^2\dd s
 \leq\tfrac12\norm{v(0)}_2^2.
\]
This does not bound its positive variation or the amplitude integral
\eqref{eq:amplitude}. The following example also excludes an automatic
viscous damping argument based just on small relative angular defect.

\begin{proposition}\label{prop:transfer}
For every $\nu>0$, $C>0$, and $0<\epsilon\leq1$, there are real
Schwartz solenoidal data $u_+$ and $u_-$ with $\norm{u_\pm}_2=1$
and the same angular mean $v$. Their defects $w_\pm=u_\pm-v$ have
equal squared $L^2$ and gradient norms, and
$\norm{w_\pm}_2=(3\epsilon/2)\norm v_2$.
Their nonlinear kinetic transfers have opposite signs, and the
unforced solution from $u_+$ satisfies at its initial time
\[
 \frac12\frac{\mathrm d}{\mathrm dt}\norm{w_+(t)}_2^2
                         >C\nu\norm{\nabla w_+}_2^2.
\]
The mean satisfies the spatial bounds
\eqref{eq:meanr}--\eqref{eq:meanz} at a time $t=-\tau_\lambda$,
with $h=1/6$ and a constant independent of the concentration parameter
$\lambda$ and of $\epsilon\in(0,1]$.
This last assertion concerns that time slice, not an entire evolution.
\end{proposition}

\begin{proof}
Set $\rho=x^2+y^2+z^2$, $r^2=x^2+y^2$, and, for $\sigma=\pm1$,
\[
 V=\curl(e^{-\rho}e_z),\qquad
 W_\sigma=\curl\big((x^2-y^2+2\sigma xy r^2)e^{-\rho}e_z\big).
\]
Both fields are horizontal, divergence free and Schwartz;
$\PP V=V$ and $\PP W_\sigma=0$. Write $g_a=(\pi/a)^{3/2}$.
Gaussian moments give
\begin{align*}
 \norm V_2^2&=2g_2,& \norm{W_\sigma}_2^2&=\tfrac92 g_2,\\
 \norm{\nabla_h W_\sigma}_2^2&=30g_2,&
 \norm{\partial_z W_\sigma}_2^2&=\tfrac92 g_2,\\
 \int W_\sigma\cdot(W_\sigma\cdot\nabla)V
       &=-\sigma\tfrac{16}{9}g_3.
\end{align*}
For example, the radial--azimuthal entry of
$\PP(W_\sigma\otimes W_\sigma)$ is
$2\sigma r^4e^{-2\rho}$, while
$\partial_r V_\theta-V_\theta/r=-4r^2e^{-\rho}$;
their product directly gives the last integral.

Let $\mu=\lambda^{2/3}$, $\lambda\geq1$, and put
\[
 C_\epsilon=g_2(2+\tfrac92\epsilon^2),\qquad
 A=\frac{\lambda\sqrt\mu}{\sqrt{C_\epsilon}},\qquad
 u_\sigma(x,y,z)=A(V+\epsilon W_\sigma)(\lambda x,\lambda y,\mu z).
\]
The absence of an axial component preserves incompressibility under
this anisotropic scaling. Orthogonality proves the stated unit norm
and relative defect norm. The positive nonlinear injection $I_+=-T_+$
and its viscous cost $\nu D_+$ obey
\[
 \frac{I_+}{\nu D_+}
 =\frac{A\lambda(16g_3/9)}
        {\nu g_2(30\lambda^2+\tfrac92\mu^2)}
 \longrightarrow\infty.
\]
Indeed this ratio grows as a positive constant times
$\lambda^{1/3}/\nu$, with $\epsilon$ fixed. Choose it larger than
$C+1$ and apply \eqref{eq:kinetictransfer}. For $\sigma=-1$ the
injection changes sign while the dissipation is unchanged. These are
jets of actual local unforced solutions from Schwartz data; no imposed
strain or truncated source is used. Continuity also gives a nonzero
initial interval with the strict inequality, but no uniform length
of that interval is asserted.

Finally take $\tau_\lambda=\lambda^{-2}$. The mean has only an
azimuthal component, whose $(a,b)$ derivative is bounded by
\[
 C_{ab}\lambda^{4/3+a+(2/3)b}
 =C_{ab}\tau_\lambda^{-1/2-1/6-a/2-(1/2-1/6)b}.
\]
The constants are uniform for $0<\epsilon\leq1$ because
$C_\epsilon\geq2g_2$. This proves the asserted spatial bounds.
\end{proof}

The example refutes zero-allowance viscous absorption of kinetic
mean--defect transfer, even with small relative defect and the stated
spatial geometry. It does not refute a data-dependent history allowance,
control of repeated interactions, or Theorems~\ref{thm:main}
and~\ref{thm:amplitude}. In particular these transfers are kinetic
energy transfers, not vorticity stretching integrals.

There is also an exact circulation test. The initial mean circulation
$\Gamma=r v_\theta$ has its global maximum at $r=1/\lambda$, $z=0$,
with value $2A/(e\lambda)$. For the base Gaussian defect,
\[
 rF_\theta=4\sigma r^4(2r^2-3)e^{-2\rho},
 \qquad F=-\PP[(W_\sigma\cdot\nabla)W_\sigma].
\]
At that maximum the radial and axial diffusion of the base circulation
are $-8/e$ and $-4/e$. Consequently the actual initial mean equation gives
\begin{equation}\label{eq:gaussiancirculation}
 \Gamma_t(1/\lambda,0)
 =-\frac{4\nu A}{e\lambda}(2\lambda^2+\mu^2)
                         -\frac{4\sigma\epsilon^2A^2}{e^2}.
\end{equation}
For $\sigma=-1$, choose $\lambda$ still larger so that
$\epsilon^2 A/(\nu\lambda)>e[2+(\mu/\lambda)^2]$.
The value at the old maximum then increases initially, and hence exceeds
the initial global maximum. Its energy comes from the defect in
\eqref{eq:kinetictransfer}. Thus the maximum principle for exactly
axisymmetric circulation cannot be extended by dropping the Reynolds
torque, even for arbitrarily small relative angular defect. This is an
initial growth example, not repeated spin-up or singularity formation.

\section{Scope of the calculation}

The theorem allows an unbounded angular fluctuation, but requires its
history to satisfy \eqref{eq:history} and separately assumes the mean
geometry. Neither hypothesis is established here for general data.
Theorem~\ref{thm:amplitude} removes a separate gradient hypothesis in
the whole-space setting, but its weighted amplitude condition remains
unproved for arbitrary data. Its localization constants may depend on
the solution on a regular annulus; they are not uniform data-only bounds.
An original force uniformly bounded in spatial $C^2$ can also be
included: its torque has finite $J_f$, its rescaled vorticity source
vanishes, and its enstrophy contribution is bounded by $C(Y+1)$.
Thus a distinction between forced and unforced flows must enter a
proof of the hypotheses, not be inferred from this criterion.

The accompanying symbolic calculation checks the mean equation,
cylindrical stress divergence, lifted force identity, and the
stretching integration by parts on explicit polynomial fields with
all six stress components nonzero. It also checks the Hardy constant
on a symbolic power profile, and the Gaussian moments and complete
convection work in Proposition~\ref{prop:transfer}. The polynomial fields
are local algebra tests, whereas the Gaussian fields are finite-energy
initial data. The program does not integrate solutions or certify the
analytic proof. No statement in this
note has been added as an axiom or a theorem to the Lean library.

\begin{thebibliography}{9}
\bibitem{CIV}
P. Constantin, M. Ignatova, and V. Vicol,
\emph{Regularity of asymptotically axisymmetric solutions to the 3D
Navier--Stokes equations with analytic forcing},
\href{https://arxiv.org/abs/2609.20803}{arXiv:2609.20803v1}, 2026.
\bibitem{GKT}
S. Gustafson, K. Kang, and T.-P. Tsai,
\emph{Interior regularity criteria for suitable weak solutions of the
Navier--Stokes equations}, Comm. Math. Phys. \textbf{273} (2007), 161--176,
\href{https://arxiv.org/abs/math/0607114}{arXiv:math/0607114}.
\end{thebibliography}
\end{document}
